\subsection{实验原理}
矩阵乘法$C_{m\times n} = A_{m \times k}\times B_{k \times n}$,要求第一个矩阵的列数等于第一个矩阵的行数,一般的实现方法为3层循环实现,
最内层最内层循环分别遍历A的行和B的列,其算法时间复杂度为$O(N^3)$.在现代计算机体系结构下,制约其性能的核心因素通常不是计算速度,而是存储层级:
\begin{itemize}
    \item 空间局部性:当访问矩阵B时将固定列而遍历行,对于按行优先存储的编程语言,这将跳跃访问内存地址,空间局部性差.
    \item 当矩阵的规模较大而超过L/L2Cache的大小时,内存MISS率将大大提高,程序执行效率低下.通过矩阵分块,使分块后的矩阵能完整装入L1/L2Cache,将减少对主存的访问.
    \item 多核并行:通过将维度任务划分,平衡CPU的核心负载,避免因负载不均导致的程序整体效率低.
\end{itemize}